% Apartado 2.3 - Filtro temporal FIR o IIR (20 puntos) - BLOQUE 3
%
% Los ocho incisos del enunciado:
%   a) banda de paso, banda de rechazo, frecuencias de corte, rizado, atenuacion
%   b) justificacion del tipo y del orden
%   c) coeficientes de la ecuacion de diferencias
%   d) funcion de transferencia H(z)
%   e) polos y ceros, region de convergencia causal, verificacion de estabilidad
%   f) respuesta en magnitud, fase y retardo de grupo
%   g) aplicacion a los tres casos, aunque el diseno se optimice para uno
%   h) cuantizacion de coeficientes, saturacion y memoria en el microcontrolador
%
% Se puede usar biblioteca para DISENAR, pero hay que aplicar la ecuacion de
% diferencias propia y demostrar que da la misma salida que la funcion de
% biblioteca dentro de una tolerancia declarada.
% Codigo: src/filtros_temporales.py   Figuras: fig_23_*.png
% Salidas: resultados/salidas/caso_N_temporal.wav para los tres casos

\section{Filtro temporal FIR o IIR}
\label{sec:temporal}

\textit{(pendiente — bloque 3)}

\subsection{Especificación y elección del filtro}

\subsection{Función de transferencia, polos y ceros}

\subsection{Respuesta en frecuencia y retardo de grupo}

\subsection{Verificación contra la implementación de biblioteca}

\subsection{Aplicación a los tres casos}

\subsection{Cuantización, saturación y memoria}
